\documentclass[10pt,a4paper]{article}
\usepackage[margin=1.8cm]{geometry}
\usepackage[T1]{fontenc}\usepackage{lmodern}\usepackage{amssymb}
\usepackage{booktabs,longtable,array,xcolor,enumitem}
\usepackage[hidelinks]{hyperref}
% Fillable form: one exclusive choice and one comment box per pair.
\newcommand{\viewfield}[1]{\ChoiceMenu[popdown,name=view_#1,width=2.1cm,height=1.2em,
  bordercolor={0.6 0.6 0.6},charsize=8pt,default= ]{}{ = ,agree=agree,disagree=disagree,unsure=unsure}}
\newcommand{\commentfield}[1]{\TextField[name=comment_#1,multiline=true,width=5.2cm,
  height=1.35cm,bordercolor={0.7 0.7 0.7},charsize=8pt]{}}
\renewcommand{\familydefault}{\sfdefault}
\newcolumntype{L}[1]{>{\raggedright\arraybackslash}p{#1}}
\setlength{\parindent}{0pt}\setlength{\parskip}{5pt}
\pagestyle{plain}
\begin{document}
\begin{Form}
{\Large\bfseries Which lymphoid cancers should learn from each other?}\par
{\large Request for review of a list fixed before testing}\par
{\footnotesize\color{black!55}Draft for the study haematologists. Fill in this PDF directly
(Adobe Reader or Preview), save it and send it to Mikkel Werling.}\par
Your name: \TextField[name=reviewer,width=7cm,height=1.1em,bordercolor={0.7 0.7 0.7}]{}

\textbf{Background.} We have built one prediction model for patients starting first-line treatment
for any of 24 lymphoid cancers (37,531 Danish patients). For each patient it predicts the risk of
69 outcomes over the following two years, from relapse and second-line treatment to infection,
cytopenia, organ toxicity and hospital admission. Because the model learns from all diseases at
once, what it learns in one disease can be used for another.

\textbf{What we will test.} After the model is trained, we can switch off what it learned from one
disease and see how much worse its predictions become for another. We want to check whether the
model relies on other diseases \emph{where clinical knowledge says it should}. To make this a fair
test, we must write down the expected pairs \textbf{before} we look at the results on the test
patients. Once the test data are opened, the list can no longer be changed. Your clinical
judgement now is therefore what makes the test meaningful.

\textbf{What we ask of you.}
\begin{enumerate}[leftmargin=1.4em,itemsep=1pt,topsep=0pt]
  \item For each pair below, tick \emph{agree}, \emph{disagree} or \emph{unsure}: would you
  expect what is learned in the first disease to help predict the stated outcomes in the second?
  Add a comment where useful. We deliberately give no reasons, so as not to steer your view.
  \item Suggest pairs we have missed. Any disease can be paired with any other, for any group of
  outcomes; the diseases and outcomes are listed at the end. Pairs that would not be obvious from
  the classification alone, such as diseases that share a treatment or a typical complication, are
  especially useful.
  \item Check the \emph{negative controls} (the last three rows): pairs we expect \emph{not} to
  inform each other. They show whether the method finds a signal where none should exist.
\end{enumerate}

{\small\setlength{\tabcolsep}{4pt}
\begin{longtable}{L{0.55cm}L{4.6cm}L{2.8cm}L{2.3cm}L{5.4cm}}
\toprule
& Pair & For predicting & Your view & Comment\\ \midrule\endfirsthead
\toprule
& Pair & For predicting & Your view & Comment\\ \midrule\endhead
\multicolumn{5}{l}{\bfseries Pairs expected to inform each other}\\ \midrule
1 & Learn from \textbf{Chronic lymphocytic leukaemia (CLL)}\newline to predict for \textbf{Small lymphocytic lymphoma (SLL)} & all outcomes & \viewfield{H1} & \commentfield{H1} \\ \midrule
2 & Learn from \textbf{Multiple myeloma (MM)}\newline to predict for \textbf{Plasma-cell leukaemia (PCL)} & disease course and later malignancy & \viewfield{H2} & \commentfield{H2} \\ \midrule
3 & Learn from \textbf{Extranodal marginal-zone lymphoma (EMZL)}\newline to predict for \textbf{Splenic marginal-zone lymphoma (SMZL)} & disease course and later malignancy & \viewfield{H3} & \commentfield{H3} \\ \midrule
4 & Learn from \textbf{Diffuse large B-cell lymphoma (DLBCL)}\newline to predict for \textbf{Transformed follicular lymphoma (tFL)} & disease course and later malignancy & \viewfield{H4} & \commentfield{H4} \\ \midrule
5 & Learn from \textbf{Diffuse large B-cell lymphoma (DLBCL)}\newline to predict for \textbf{Richter transformation (RT)} & disease course and later malignancy & \viewfield{H5} & \commentfield{H5} \\ \midrule
6 & Learn from \textbf{Chronic lymphocytic leukaemia (CLL)}\newline to predict for \textbf{Richter transformation (RT)} & infection and antimicrobial use & \viewfield{H6} & \commentfield{H6} \\ \midrule
7 & Learn from \textbf{Multiple myeloma (MM)}\newline to predict for \textbf{Chronic lymphocytic leukaemia (CLL)} & infection and antimicrobial use & \viewfield{H7} & \commentfield{H7} \\ \midrule
8 & Learn from \textbf{Chronic lymphocytic leukaemia (CLL)}\newline to predict for \textbf{Multiple myeloma (MM)} & infection and antimicrobial use & \viewfield{H8} & \commentfield{H8} \\ \midrule
9 & Learn from \textbf{Multiple myeloma (MM)}\newline to predict for \textbf{Amyloidosis} & renal injury and support & \viewfield{H9} & \commentfield{H9} \\ \midrule
10 & Learn from \textbf{Burkitt lymphoma (BL)}\newline to predict for \textbf{Lymphoblastic lymphoma (LBL)} & neutrophil, white-cell and platelet findings & \viewfield{H10} & \commentfield{H10} \\ \midrule
11 & Learn from \textbf{Burkitt lymphoma (BL)}\newline to predict for \textbf{Lymphoblastic lymphoma (LBL)} & metabolic and cell-turnover markers & \viewfield{H11} & \commentfield{H11} \\ \midrule
12 & Learn from \textbf{Peripheral T-cell lymphoma, NOS (PTCL)}\newline to predict for \textbf{T-cell lymphoma, unspecified (TCL)} & disease course and later malignancy & \viewfield{H12} & \commentfield{H12} \\ \midrule
13 & Learn from \textbf{Chronic lymphocytic leukaemia (CLL)}\newline to predict for \textbf{Lymphoplasmacytic lymphoma (LPL)} & cardiovascular and thrombotic events & \viewfield{H13} & \commentfield{H13} \\ \midrule
\multicolumn{5}{l}{\bfseries Negative controls: pairs expected \emph{not} to inform each other}\\ \midrule
C1 & Learn from \textbf{Amyloidosis}\newline to predict for \textbf{Hodgkin lymphoma (HL)} & all outcomes & \viewfield{C1} & \commentfield{C1} \\ \midrule
C2 & Learn from \textbf{Anaplastic large-cell lymphoma (ALCL)}\newline to predict for \textbf{Solitary plasmacytoma} & all outcomes & \viewfield{C2} & \commentfield{C2} \\ \midrule
C3 & Learn from \textbf{Transformed follicular lymphoma (tFL)}\newline to predict for \textbf{Hairy-cell leukaemia (HCL)} & disease course and later malignancy & \viewfield{C3} & \commentfield{C3} \\ \midrule
\end{longtable}}

\textbf{Pairs we have missed} (learn from which disease, to predict which outcomes in which
disease, and why):\par
\TextField[name=missed_pairs,multiline=true,width=\linewidth,height=3.2cm,
  bordercolor={0.7 0.7 0.7},charsize=9pt]{}\par\medskip
\textbf{Other comments:}\par
\TextField[name=other_comments,multiline=true,width=\linewidth,height=2.4cm,
  bordercolor={0.7 0.7 0.7},charsize=9pt]{}

\newpage
{\large\bfseries The 24 diseases}\par
\begin{itemize}[leftmargin=0pt,itemsep=2pt,topsep=2pt]
\item[] \textbf{Large B-cell lymphomas}: Diffuse large B-cell lymphoma (DLBCL); Richter transformation (RT); Transformed follicular lymphoma (tFL)
\item[] \textbf{Burkitt and lymphoblastic lymphomas}: Burkitt lymphoma (BL); Lymphoblastic lymphoma (LBL)
\item[] \textbf{Mantle-cell lymphoma}: Mantle-cell lymphoma (MCL)
\item[] \textbf{Indolent B-cell lymphomas}: Follicular lymphoma (FL); Lymphoplasmacytic lymphoma (LPL); Extranodal marginal-zone lymphoma (EMZL); Splenic marginal-zone lymphoma (SMZL); Nodal marginal-zone lymphoma (NMZL); B-cell lymphoma, unspecified (BCL)
\item[] \textbf{CLL and SLL}: Chronic lymphocytic leukaemia (CLL); Small lymphocytic lymphoma (SLL)
\item[] \textbf{Hairy-cell leukaemia}: Hairy-cell leukaemia (HCL)
\item[] \textbf{Hodgkin lymphoma}: Hodgkin lymphoma (HL)
\item[] \textbf{T-cell lymphomas}: Peripheral T-cell lymphoma, NOS (PTCL); Angioimmunoblastic T-cell lymphoma (AITL); Anaplastic large-cell lymphoma (ALCL); T-cell lymphoma, unspecified (TCL)
\item[] \textbf{Plasma-cell neoplasms}: Multiple myeloma (MM); Plasma-cell leukaemia (PCL); Solitary plasmacytoma
\item[] \textbf{Amyloidosis}: AL amyloidosis
\end{itemize}

{\large\bfseries The 69 outcomes, by clinical domain}\par
Each outcome is predicted at 30, 60 and 90 days and at 6, 12 and 24 months after the start of
first-line treatment. Laboratory grades are based on the CTCAE laboratory thresholds; pairs can name a
whole domain or ``all outcomes''.
\begin{itemize}[leftmargin=0pt,itemsep=2pt,topsep=2pt]
\item[] \textbf{Disease course and later malignancy}: overall survival; second-line treatment; treatment failure; treatment failure or transformation; new haematological malignancy; new malignancy (any)
\item[] \textbf{Infection and antimicrobial use}: sepsis; serious infection (including invasive fungal infection and pneumocystosis); antibiotic treatment; intravenous antimicrobial use; prescription antimicrobial use; infection diagnosis or antimicrobial use
\item[] \textbf{Erythroid findings and transfusion}: anaemia grade 2 and 3; red-cell transfusion; any transfusion; haptoglobin decreased
\item[] \textbf{Neutrophil, white-cell and platelet findings}: neutropenia grade 2 and 3; leukopenia grade 2 and 3; thrombocytopenia grade 2 and 3; platelet transfusion; severe cytopenia
\item[] \textbf{Lymphocyte and eosinophil counts}: lymphocytosis grade 2 and 3; eosinophilia
\item[] \textbf{Renal injury and support}: acute kidney injury; creatinine increased grade 2 and 3; creatinine rise above 1.5 times the pre-treatment value; dialysis
\item[] \textbf{Liver injury and coagulation}: ALT increased grade 2 and 3; alkaline phosphatase elevation, two levels; bilirubin increased grade 2 and 3; hepatotoxicity; INR increased grade 2 and 3
\item[] \textbf{Sodium and potassium}: hyponatraemia grade 2; sodium below 125 mmol/L; potassium below 3.0 mmol/L; hypernatraemia grade 2 and 3; hyperkalaemia grade 2 and 3
\item[] \textbf{Calcium and magnesium}: hypercalcaemia grade 2 and 3; hypocalcaemia grade 2 and 3; hypomagnesaemia grade 2
\item[] \textbf{Metabolic and cell-turnover markers}: hypertriglyceridaemia grade 2 and 3; hypoalbuminaemia grade 2 and 3; hyperphosphataemia; hyperuricaemia; LDH increased
\item[] \textbf{Cardiovascular and thrombotic events}: atrial fibrillation; heart failure; stroke or transient ischaemic attack; venous thromboembolism
\item[] \textbf{Acute care and other events}: acute hospitalisation; emergency department stay; respiratory support; bone fracture
\end{itemize}
\end{Form}
\end{document}
